Does deep learning improve forecasts of Value-at-Risk and Expected Shortfall when it is given more information, or only when it is given a better architecture? We study pre-open forecasts for 24 equity indices, using the returns of foreign markets that close before each market opens. All models are trained or evaluated with the strictly consistent FZ0 score. A pre-registered protocol develops the models on twelve indices and tests the frozen models exactly once on twelve further indices. Cross-market information is valuable: adding two overnight-news terms to a GJR-GARCH model lowers the loss at every risk level (Diebold--Mariano $t\approx-3.9$ in both samples), with the largest gains in Asian markets after U.S. and European sell-offs. Neural networks, however, do not extract this information. An attention-based contagion network, a neural news-impact GARCH with a proven stability property, and versions trained with an efficiency-improving composite score all fail to beat the two-parameter linear model. An ablation shows that the contagion network gains essentially nothing from the news, which the linear model turns into a large improvement. Forecast combinations that include both families are the most accurate in both samples and spend the fewest days in the Basel yellow zone. In tail-risk forecasting from daily data, information and inductive bias matter more than architecture.
